\newpage
\section{Conclusion}

This thesis investigated whether existing New York City traffic-camera infrastructure 
could be used as a practical source of vehicle composition data for traffic and 
emissions analysis. To address this question, a low-resolution urban traffic dataset 
was constructed from 511NY and NYSDOT camera imagery, annotated using a MOVES-aligned 
6-class vehicle taxonomy, and used to evaluate several modern object-detection 
architectures. Across the models tested, the YOLO family performed best for this 
setting, with the baseline YOLOv11s model achieving an AP50 of 0.804 and an mAP of 
0.615 on the validation set. After targeted refinement, the final selected YOLOv11s 
configuration reached an AP50 of 0.857 and an mAP of 0.661, showing that careful 
data curation and hard negative mining were more effective in this case than simply 
changing to a more complex architecture or adding more untargeted data.

The results also showed that the design of the taxonomy and the treatment of ambiguous 
annotations were central to the success of the pipeline. By grouping visually similar 
vehicle body styles into broader MOVES-relevant classes and treating 
\verb|unknown_vehicle| regions as ignore areas rather than a learnable class, the 
detector was able to focus on vehicle distinctions that were both visually reliable 
and meaningful for downstream environmental analysis. This made the final system a 
better match for the limits of 352$\times$240 traffic imagery, where many vehicles 
are too distant, too small, or too occluded to support fine-grained labeling. In 
that sense, the strongest gains in this study came not only from model selection, 
but from aligning the dataset and class structure with the actual information present 
in the images.

Beyond validation performance, the experiments provided a clearer picture of where 
the detector generalized well and where important limitations remained. On the 
held-out test set, the final model achieved an overall AP50 of 0.772 and an mAP of 
0.607, indicating that the selected improvements transferred beyond the validation 
split. On the separate cross-camera benchmark of unseen locations, performance fell 
to an AP50 of 0.658 and an mAP of 0.488, showing that transfer to new camera views 
was still more difficult. The comparison between real 2025 NYC ATR counts further 
showed that the pipeline was better at recovering the hourly pattern of traffic than exact count magnitude. 
For example, Atlantic Avenue at Boerum Place produced the closest count agreement, with 
Pearson correlations of 0.844 and 0.803 and MAPE values of 22.96\% and 27.28\%, 
while Pelham Parkway maintained strong hourly-profile agreement, with correlations 
up to 0.906, but still severely undercounted total vehicles. These findings suggest 
that the current system is already useful for broad fleet characterization and traffic 
pattern analysis, but that exact count recovery and camera-to-camera transfer remain 
open challenges.

The results support the overall feasibility of using public traffic 
camera imagery for intersection-level vehicle classification in support of transportation 
and environmental analysis. At the same time, the work also shows that this approach 
is best understood as a first stage rather than a complete end-to-end emissions-monitoring 
solution. The current pipeline can estimate vehicle composition within snapshots and 
capture broad temporal traffic structure, but extending it to fully calibrated traffic 
volumes and MOVES-ready emissions inputs will require broader training data, stronger 
cross-camera adaptation, temporal tracking, and additional operational information such 
as vehicle speed. Even with these limitations, the study demonstrates that existing 
traffic-camera networks can provide useful transportation intelligence without requiring 
new physical sensing infrastructure. In that sense, this work establishes a practical 
foundation for future camera-based traffic and emissions monitoring in urban environments.

